\section{渐开线齿轮}
\markboth{49. 渐开线齿轮}{49. 渐开线齿轮}

常见齿轮的牙齿，侧面轮廓大多做成圆的渐伸线。工厂里把圆的渐伸线叫做圆的渐开线，所以这种齿轮叫做渐开线齿轮。

\begin{wrapfigure}[19]{l}{0.46\textwidth}
  \centering
  \vspace{-4pt}
  \begin{tikzpicture}[scale=0.72, >=Stealth, line cap=round, line join=round]
    % z=8 宽大饱满渐开线齿轮，alpha=25°, ra=2.12 严格满足 2*rp - ra > rb，两齿在 M 精确外切零交叉干涉
    \def\rp{1.95}
    \def\alphaP{25}
    \pgfmathsetmacro{\rb}{\rp*cos(\alphaP)}
    \def\rf{1.70}
    \def\ra{2.12}
    \def\wb{12.968} % z=8: 11.25 + inv(25 deg) in degrees
    \pgfmathsetmacro{\tmax}{sqrt((\ra/\rb)^2 - 1)} % rad ~ 0.6626
    \def\sM{0.28} % 接触点 M 距节点 B 沿公切线左上方向的距离

    \coordinate (O1) at (0, \rp);
    \coordinate (O2) at (0, -\rp);
    \coordinate (B)  at (0, 0);
    \coordinate (M)  at ({-\sM*cos(\alphaP)}, {\sM*sin(\alphaP)});

    % 连心线 O1 O2 与旋转中心十字线
    \draw[densely dash dot] (0, {\rp+0.34}) -- (0, {-\rp-0.34});
    \draw[densely dash dot] (-0.60, \rp) -- (0.60, \rp);
    \draw[densely dash dot] (-0.60, -\rp) -- (0.60, -\rp);

    % 两个基圆弧（虚线）
    \draw[densely dashed] ($(O1)+(222:\rb)$) arc (222:318:\rb);
    \draw[densely dashed] ($(O2)+(42:\rb)$) arc (42:138:\rb);

    % 公切线 MN（向右下充分延长出齿轮 II 外缘，确保 N 位于纯空白区）
    \draw[thick] (-1.58, {1.58*tan(\alphaP)}) -- (1.88, {-1.88*tan(\alphaP)});
    \node[above right=-1pt, font=\scriptsize] at (1.86, {-1.86*tan(\alphaP)}) {$N$};

    % 齿轮 I 齿廓 (z=8, pitch=45 deg；第 2 颗牙右廓面在 M 点与齿轮 II 精确外切无干涉)
    \begin{scope}[shift={(O1)}]
      \foreach \ang [count=\i] in {204.67, 249.67, 294.67, 339.67} {
        \pgfmathsetmacro{\thu}{\ang+\wb}
        \pgfmathsetmacro{\thl}{\ang-\wb}
        \draw[thick, curveblue]
          (\thl:\rf) -- (\thl:\rb)
          -- plot[domain=0:\tmax, samples=14, variable=\t]
            ({\rb*(cos(\thl + \t*180/pi) + \t*sin(\thl + \t*180/pi))},
             {\rb*(sin(\thl + \t*180/pi) - \t*cos(\thl + \t*180/pi))})
          -- plot[domain=\tmax:0, samples=14, variable=\t]
            ({\rb*(cos(\thu - \t*180/pi) - \t*sin(\thu - \t*180/pi))},
             {\rb*(sin(\thu - \t*180/pi) + \t*cos(\thu - \t*180/pi))})
          -- (\thu:\rf);
        \ifnum\i<4
          \pgfmathsetmacro{\nextl}{\ang + 45.0 - \wb}
          \draw[thick, curveblue] (\thu:\rf) arc (\thu:\nextl:\rf);
        \fi
      }
    \end{scope}

    % 齿轮 II 齿廓 (z=8, pitch=45 deg；第 2 颗牙左廓面在 M 点与齿轮 I 精确外切无干涉)
    \begin{scope}[shift={(O2)}]
      \foreach \ang [count=\i] in {42.82, 87.82, 132.82} {
        \pgfmathsetmacro{\thu}{\ang+\wb}
        \pgfmathsetmacro{\thl}{\ang-\wb}
        \draw[thick, curveblue]
          (\thl:\rf) -- (\thl:\rb)
          -- plot[domain=0:\tmax, samples=14, variable=\t]
            ({\rb*(cos(\thl + \t*180/pi) + \t*sin(\thl + \t*180/pi))},
             {\rb*(sin(\thl + \t*180/pi) - \t*cos(\thl + \t*180/pi))})
          -- plot[domain=\tmax:0, samples=14, variable=\t]
            ({\rb*(cos(\thu - \t*180/pi) - \t*sin(\thu - \t*180/pi))},
             {\rb*(sin(\thu - \t*180/pi) + \t*cos(\thu - \t*180/pi))})
          -- (\thu:\rf);
        \ifnum\i<3
          \pgfmathsetmacro{\nextl}{\ang + 45.0 - \wb}
          \draw[thick, curveblue] (\thu:\rf) arc (\thu:\nextl:\rf);
        \fi
      }
    \end{scope}

    % 旋转方向箭头
    \draw[->, thick] ($(O1)+(208:1.26)$) arc (208:250:1.26);
    \draw[->, densely dashed, thick] ($(O2)+(152:1.26)$) arc (152:110:1.26);

    % 标注（接触点 M 与节点 B 均带清晰实心点并引至无线条空白区）
    \fill (O1) circle (1.4pt) node[above right=-1pt, font=\scriptsize] {$O_1$};
    \fill (O2) circle (1.4pt) node[below right=-1pt, font=\scriptsize] {$O_2$};
    \fill (B)  circle (1.5pt);
    \draw[thin, black!70] (B) -- (0.34, 0.30) node[above right=-2pt, font=\scriptsize] {$B$};
    \fill (M)  circle (1.5pt);
    \draw[thin, black!70] (M) -- (-0.54, 0.50) node[left=-1pt, font=\scriptsize] {$M$};
    \node[right, font=\scriptsize] at (0.14, {\rp-0.62}) {齿轮 I};
    \node[right, font=\scriptsize] at (0.14, {-\rp+0.62}) {齿轮 II};
  \end{tikzpicture}
  \caption*{图\ 49.1}
\end{wrapfigure}

图 49.1 所示为一对渐开线齿轮啮合的情形。点 $O_1$ 和 $O_2$ 分别是这两个齿轮的旋转中心，虚线画的圆是基圆。齿轮 I 每个牙齿两侧的轮廓曲线都是圆 $O_1$ 的渐伸线，齿轮 II 牙齿两侧的轮廓曲线是圆 $O_2$ 的渐伸线。设在某一瞬间这两个齿轮在点 $M$ 接触，那么它们的齿廓曲线在点 $M$ 有公共的切线，因而也有公共的法线 $MN$。根据圆的渐伸线的法线的性质（第 48 节），$MN$ 必为基圆 $O_1$ 和 $O_2$ 的公切线，因而是一条定直线。设 $MN$ 与两圆连心线 $O_1O_2$ 的交点为 $B$，则 $B$ 是定点。

当主动齿轮 I 旋转时，在接触点 $M$ 沿齿廓法线 $MN$ 的方向推动齿轮 II，使齿轮 II 跟着它旋转。在旋转的过程中，接触点 $M$ 的位置不断变化，但是齿轮 I 对齿轮 II 的推力却始终通过定点 $B$。这种受力情况，相当于摩擦轮传动：假想有一对橡皮轮，一个的圆心是 $O_1$、半径是 $O_1B$，另一个的圆心是 $O_2$、半径是 $O_2B$，那么当第一个橡皮轮转动时，由于摩擦力的作用，在 $B$ 点处始终有一个推力作用到第二个橡皮轮上，使它跟随着转动。这两个转动的橡皮轮，在接触点 $B$ 的线速度相同，因而旋转的角速度总是与它们的半径成反比。由此推出渐开线齿轮 I 和 II 旋转的角速度之比，也总是等于 $O_2B:O_1B$，这是一个定值，与时间无关。这样就能保证当齿轮 I 匀速旋转时，齿轮 II 的旋转也是匀速的。

主动齿轮旋转的角速度与从动齿轮旋转的角速度的比，叫做传动比。如果齿轮的齿廓曲线没有经过特殊选择，一般说来，在不同瞬间的传动比也是不同的。但是如果选择圆的渐伸线作齿廓曲线，就能得到固定的传动比，使传动过程平稳，并且能得到具有预期的角速度的匀速旋转运动。
